\section{Affine families: tractable global queries and a geometric extraction step}
\label{sec:affine-families-review}

Commit \texttt{58ee11a97d} supplies
\path{docs/incoming/unknot_affine_families_20261008.zip}. All 466 manifest
entries match the original archive. The delivery combines three separate
interfaces: optimizing a supplied family of cyclic translations, extracting
partial interval maps from actual normal coordinates, and inheriting scalar
endomorphism algebras through a verified direct-sum decomposition. The first
two are complementary but are not yet connected by a complete geometric
adapter. The third overlaps our earlier scalar-locality and static-corner
work, with additional transport and one-generator optimizations worth
measuring separately. This review preserves these distinctions.

\subsection{A factor-free optimum over an affine modular family}

Let $R=\mathbb Z/m\mathbb Z$, and let feasible parameters satisfy $Az=b$.
For a connected base whose monodromies are translations by the entries of
$h(z)=c+Hz$, the number of covering components is
\[
 \kappa(z)=\gcd(m,h_1(z),\ldots,h_\beta(z)).
\]
Solving the seam congruences gives the full affine coset $z=z_0+Bt$.
Its holonomy image is $v+Ut$, with $v=c+Hz_0$ and $U=HB$. The minimum
component count is exactly the gcd of $m$ and all entries of $v$ and $U$.
This is stronger than a lower bound: some assignment attains that gcd.

The constructive step merges two vectors $a,b$. Divide them and $m$ by
their common content $g$, and put $M=m/g$ and
$d=\gcd(M,a_1/g,\ldots,a_\beta/g)$. If $d=1$, retain $a$. Otherwise
iterate $D\leftarrow\gcd(M,D^2)$ from $D=d$ until stable, and use
$a+(M/D)b$. At termination $D$ contains the full prime-power parts of
$M$ supported on $d$. On those primes $a/g$ vanishes and $(M/D)b/g$
has a nonzero coordinate; on the other primes the added term vanishes
and $a/g$ has a nonzero coordinate. The new vector has content $g$.
This proof uses prime valuations to justify the algorithm, while the
algorithm itself only computes gcds. Saturation doubles positive
valuations, giving $O(\log\log(M+2))$ iterations. For example $m=12$,
$a=2$, $b=1$ needs the saturated coefficient three; the unsaturated
coefficient six would give content four rather than one.

A row-by-row modular solver maintains the complete affine coset using
unimodular two-column gcd transformations and a single congruence restriction
per row. All stored residues are reduced modulo $m$. The parameterization
may be redundant, but uniform parameters have equal-size fibres, so dividing
the assignment count by the size of each new scalar image counts distinct
feasible assignments correctly. Explicit matrix dimensions and witness
lengths are charged; a binary declaration of an enormous number of unused
variables is not treated as a constant-size explicit-variable input.

For a proposed minimum $\mu\mid m$, set $s=m/\mu$. A feasible witness
$z_*$ with $\gcd(m,c+Hz_*)=\mu$, together with
\[
 YA=sH,\qquad Yb+sc=0\pmod m,
\]
proves optimality without a kernel solver: every feasible holonomy is
annihilated by $s$ and hence divisible by $\mu$, while $z_*$ attains it.
An infeasibility witness is a row $y$ with $yA=0$ and $yb\ne0$ modulo $m$.
The finite-module annihilator identity guarantees existence also for
composite moduli. The implemented verifier checks these elementary
identities and does not invoke the optimizer. Its minimal certificate
contract deliberately excludes the separately computed feasible-assignment
count.

\subsection{Where the family query stops}

The surface-family compiler permits disconnected local cyclic covers,
full-boundary seams, shared affine parameters and a connected patch graph.
A spanning tree transports sheet labels; local generators and non-tree
seam voltages generate the global translation action. Seam compatibility
is the equality of peripheral translations with the specified base-circle
sign. A negative base-circle degree is not a reflection of the cyclic
fibre. The compiler recomputes the holonomy data from the original surface
schema before checking an arithmetic certificate.

The polynomial optimum does not enumerate every marked covering or preserve
all information needed for future attachments. Even connected covers can
have different genus: over a five-sheet pair of pants, peripheral
translations $(1,1,-2)$ give three boundary components and genus two,
whereas $(1,-1,0)$ give seven boundary components and genus zero. Both
covers are connected and have Euler characteristic $-5$.

The report gives a sharp warning about correlated peripheral queries.
For a graph, assign one parameter $z_v$ modulo three to each vertex and
one boundary translation $z_u-z_v$ to each edge, on a planar surface with
one additional boundary. Requiring every designated edge-boundary preimage
to be connected is exactly proper three-colouring. This proves NP-completeness
for that unrestricted affine peripheral query, not for unknot recognition,
and does not unconditionally exclude quasi-polynomial time. Similarly,
counting connected assignments in the family $h(z)=z$ includes computing
$\varphi(m)$; for a semiprime this recovers its factors. Neither issue
contradicts the polynomial minimum-and-witness theorem.

\subsection{From actual normal coordinates to local cut data}

The normal extractor accepts reciprocal tetrahedral face pairings and an
admissible seven-coordinate normal vector. It checks nonnegative coordinates,
quadrilateral compatibility and matching equations. The ambient manifold
condition remains a precondition; no PD-to-exterior conversion is supplied.
Each tetrahedron has at most four triangle stacks and one quadrilateral
stack. Face arcs are ordered from the cut-off vertex, with triangles before
quadrilaterals. Quadrilateral order reverses on the opposite side of its
chosen vertex partition. Equal face ranks therefore give affine maps
$i\mapsto\varepsilon i+c$ between intervals of disc indices.

There are at most five such surface bands per paired face: for $k\le3$
nonempty arc types, either side has at most $k+1$ blocks, and overlaying
the two partitions gives at most $2(k+1)-k=k+2\le5$ intersections.
Consecutive-disc prism gaps require both endpoint discs to lie in the band.
Their map is $i\mapsto i+c$ for a translation, but
$i\mapsto-i+c-1$ for a reflection. The minus one is necessary because a
gap is indexed by its smaller disc index.

If a tetrahedron has $D$ normal discs and $k$ nonempty types, it has
$D-k$ between-disc prisms and $k+1\le6$ remaining local cells. Across a
paired face at most two exceptional prism sides meet residual cells: one
in each direction where an arc block crosses a triangle--quadrilateral
transition. These are singleton gap contacts. The output size and arithmetic
work are linear in tetrahedra and polynomial in coordinate bit length,
with no multiplicity expansion.

Concatenating disc-stack intervals produces an input for interval-orbit
counting. A signed orientation double cover similarly distinguishes
orientable and nonorientable component counts once those orbits are known.
The delivery prepares these inputs; it does not implement unrestricted
Agol--Hass--Thurston reduction, all weighted component data, residual edge
attachments, or a full cut-manifold handle structure. Its partial signed
interval maps are not automatically total translations on one cyclic fibre.
A sound extraction and coverage theorem is still needed to connect them
to the affine-family optimizer.
The later integration in Section~\ref{sec:native-interval-orbits} supplies
unweighted AHT counting and aggregate topology of supplied normal surfaces;
weighted extraction and attachment-aware cutting remain open implementation
obligations.

\subsection{Scalar inheritance relative to the maintained implementation}

For a complete scalar commutant $A$ and an actual typed direct-sum child
with coordinate selector $e$, restriction identifies its full algebra with
$e(S^{-1}AS)e$: every child endomorphism extends by zero. The map from the
whole algebra to the corner is only linear in general, not multiplicative.
Canonical least-pivot reduction recovers the same ordered basis as a fresh
commutator solve, preserving bounded candidate order. Our maintained
\path{corner_split.py} already implements this exact static principle;
its reuse option remains separate from runtime defaults.

The delivery additionally selects direct transport or a precomputed packed
linear map on matrix variables. When a generator $Q$ saturates the full
commutant dimension, $A=\mathbb F_2[Q]$. A child of this commutative algebra
is generated by the restricted $Q$, so its full algebra can be recovered
from one transported generator. A primary minimal polynomial then proves
scalar locality. Our earlier dimension-saturation stopping already checks
the locality implication; the new contribution to evaluate is propagation
through one generator and its transport cost.

The delivery reports fewer equations and solves on attached-copy fixtures,
and about twofold improvements on selected saturated two-field fixtures.
Those are graded algebraic complexes without a claimed realization as
knot-scan prefixes. Its six natural-knot timings range from 0.961 to 1.022
in paired baseline/new ratios, providing no convincing recognition gain.
These supplied timings are not measurements of our current implementation;
its baseline predates several maintained changes. A blanket patch would
also overwrite our independently integrated scalar options. We therefore
retain the proof and reproduction evidence, and defer runtime promotion
until a direct comparison with the maintained backend establishes a benefit.

\subsection{Reproduction and integration status}

We reran all 65 targeted methods in the distributed source tree: 13 arithmetic,
19 family-compiler, eight normal-extractor, 13 scalar-inheritance and 12
primary-split tests. All passed. The rational oracle reproduced 1,537
geometric cases, 9,771 disc identifications and 4,850 prism identifications,
including 591 reflected prism bands. The separate Regina audit reproduced
807 inputs and 1,537 component records, including 92 nonorientable and 274
closed components; it compared full coordinate multisets, Euler characteristic,
orientability and boundary counts. Family example and peripheral-reduction
checks also passed. Logs and JSON evidence are stored under
\path{data/affine-families-*}.

These are reproduced delivery tests, not a claim that the entire distributed
source tree replaces our maintained implementation. We did not rerun its
790-test historical suite or its comparative timing experiments. The current
775-test maintained suite validates the separate braid change in
Section~\ref{sec:linear-braid-descent}. The archive also repeats the earlier
single-submission worker-communication repair; it is distinct from all three
new mathematical interfaces and should be assessed as a focused resource
handling change.

The most consequential new geometric result is the actual normal-coordinate
extractor. The next integration target is a checked native consumer of those
intervals, with general weighted orbit reduction and full residual attachments.
A polynomial family query can remove enumeration only when its invariant is
all that the caller needs. Neither component minimization nor inherited
scalar locality supplies the missing global state-count, reset or completeness
bound for the unknot algorithm.
